\documentclass[11pt,a4paper]{article}

\usepackage[margin=2.5cm]{geometry}
\usepackage{amsmath,amssymb,amsthm,mathtools}
\usepackage{booktabs,graphicx,enumitem}
\usepackage[labelfont=bf,font=small,labelsep=period]{caption}
\usepackage{algorithm}
\usepackage{algpseudocode}
\usepackage[expansion=false]{microtype}
\usepackage{placeins}
\usepackage[colorlinks=true,linkcolor=black,citecolor=black,urlcolor=black]{hyperref}
\usepackage[capitalise,noabbrev]{cleveref}

\setlength{\parindent}{1.5em}
\setlength{\parskip}{0pt}
\renewcommand{\arraystretch}{1.15}
\numberwithin{equation}{section}
\setlist[itemize]{itemsep=2pt,topsep=4pt}
\setlist[enumerate]{itemsep=2pt,topsep=4pt}

% ---------- theorem environments ----------
\theoremstyle{plain}
\newtheorem{theorem}{Theorem}[section]
\newtheorem{proposition}[theorem]{Proposition}
\newtheorem{lemma}[theorem]{Lemma}
\theoremstyle{definition}
\newtheorem{definition}[theorem]{Definition}
\theoremstyle{remark}
\newtheorem{remark}[theorem]{Remark}
\newtheorem{example}[theorem]{Example}

% ---------- notation ----------
\newcommand{\R}{\mathbb{R}}
\newcommand{\E}{\mathbb{E}}
\newcommand{\T}{^{\top}}
\newcommand{\dd}{\mathrm{d}}
\DeclareMathOperator{\Var}{Var}
\DeclareMathOperator{\KL}{KL}
\DeclareMathOperator{\Hess}{Hess}
\DeclareMathOperator{\grad}{grad}
\DeclareMathOperator{\diag}{diag}
\DeclareMathOperator*{\argmax}{arg\,max}
\newcommand{\aref}{a_{\mathrm{ref}}}
\newcommand{\Ao}{A^{\circ}}
\newcommand{\pio}{\pi_{\circ}}
\newcommand{\Ind}{\mathbf{1}}
\newcommand{\mani}{\mathcal{M}}

% ---------- figures: placeholder if the file is missing ----------
\newcommand{\figslot}[2][0.92]{%
  \IfFileExists{figures/#2.pdf}%
    {\includegraphics[width=#1\linewidth]{figures/#2.pdf}}%
    {\fbox{\parbox[c][5.5cm][c]{#1\linewidth}{\centering\small [placeholder: \texttt{\detokenize{figures/#2.pdf}}]}}}}

% ---------- author block (edit here) ----------
\newcommand{\paperauthor}{Hausel-Vargha Emil}
\newcommand{\paperaffil}{Course on Differential Geometry}

\title{The Natural Policy Gradient through Differential Geometry:\\
Theory, Algorithm and Experiments}
\author{\paperauthor\\[2pt]\small\paperaffil}
\date{October 2026}

\begin{document}
\maketitle

\begin{abstract}
\noindent
The set of policies of a Markov decision process is endowed with the structure of a manifold, and
parametrizations of policies are interpreted as charts. In this setting the differential of the
expected return is a covector, an update of the parameters is a vector, and a metric is required to
convert the former into the latter. We take as metric the Fisher metric, defined as the Hessian of the
Kullback--Leibler divergence at its minimum, and show that a step of fixed Fisher length is, to
second order, a step of fixed Kullback--Leibler size in every parametrization. In Part~I we construct
the natural policy gradient on the manifold of tabular policies in a genuine chart, derive its
closed form, prove that it is the solution of a least-squares regression of the advantage on the score
(Kakade's theorem), and describe an estimator based on a statewise regularised regression.
In Part~II we test the resulting algorithm on a Gaussian bandit and on a gridworld with bonus cells and pits.
\end{abstract}
\clearpage
\tableofcontents
\clearpage

% =====================================================================
\part{Theory: from the policy gradient to the natural gradient}
% =====================================================================

\section{Motivation}\label{sec:motivation}

\subsection{Bandits and Markov decision processes}

A multi-armed bandit has a finite set of arms $A=\{1,\dots,K\}$. Pulling the arm $a$ yields a random
reward of mean $\mu_a$. A policy is a probability vector $\pi=(\pi_a)_{a\in A}$, the arm $a$ being
chosen with probability $\pi_a$, and its quality is the expected reward
\[
  \eta(\pi)=\sum_{a\in A}\pi_a\,\mu_a .
\]

A Markov decision process (MDP) generalises this to sequences of decisions. It consists of a finite
set $S$ of non-terminal states (together with terminal states), a finite set $A$ of actions,
transition probabilities $P(s'\mid s,a)$, a reward $r_t$ received at each step, a discount factor
$\gamma\in[0,1)$, and a start state $s_0$. A policy $\pi$ assigns to every state $s\in S$ a
probability distribution $\pi(\cdot\mid s)$ on $A$. An episode $(s_0,a_0,r_0,s_1,a_1,r_1,\dots)$ is
generated by drawing $a_t\sim\pi(\cdot\mid s_t)$ and $s_{t+1}\sim P(\cdot\mid s_t,a_t)$ until a terminal
state is reached.

\begin{definition}[Gain and value functions]\label{def:gain}
The gain (discounted return) from time $t$ is the random variable
\begin{equation}\label{eq:gain}
  G_t=\sum_{k\ge 0}\gamma^{k}\,r_{t+k},
\end{equation}
with $r_u=0$ after termination. The objective, the state value, the action value and the advantage
of a policy $\pi$ are
\begin{align}
  \eta(\pi)&=\E_\pi[G_0], &
  V^\pi(s)&=\E_\pi[G_t\mid s_t=s],\nonumber\\
  Q^\pi(s,a)&=\E_\pi[G_t\mid s_t=s,\,a_t=a], &
  A^\pi(s,a)&=Q^\pi(s,a)-V^\pi(s), \label{eq:values}
\end{align}
where $\E_\pi$ denotes the expectation over the actions drawn from $\pi$ and over the transitions.
\end{definition}

By construction
\[
  V^\pi(s)=\sum_{a}\pi(a\mid s)\,Q^\pi(s,a),
  \qquad
  \sum_{a}\pi(a\mid s)\,A^\pi(s,a)=0,
  \qquad
  \eta(\pi)=V^\pi(s_0).
\]
A bandit is the MDP with a single state and episodes of one step; then $Q(a)=\mu_a$, $V=\eta$ and
$A_a=\mu_a-\eta$.

\begin{definition}[Discounted visitation]\label{def:visitation}
Let $T$ denote the termination time. The discounted visitation of the state $s$, its total mass and
its normalisation are
\begin{equation}\label{eq:visitation}
  u^\pi(s)=\sum_{t\ge 0}\gamma^{t}\,\Pr\nolimits_\pi\!\bigl(s_t=s,\ t<T\bigr),
  \qquad
  U^\pi=\sum_{s}u^\pi(s),
  \qquad
  d^\pi=\frac{u^\pi}{U^\pi},
\end{equation}
and $\rho^\pi(s,a)=d^\pi(s)\,\pi(a\mid s)$ is the discounted state--action distribution.
\end{definition}

\begin{lemma}\label{lem:visitation}
For every function $f$ of state--action pairs,
\begin{equation}\label{eq:visitation-identity}
  \E_\pi\Bigl[\sum_{t<T}\gamma^{t}f(s_t,a_t)\Bigr]
  =\sum_{s}u^\pi(s)\sum_{a}\pi(a\mid s)f(s,a)
  =U^\pi\,\E_{(s,a)\sim\rho^\pi}\bigl[f(s,a)\bigr].
\end{equation}
\end{lemma}

\begin{proof}
Since $a_t$ is drawn from $\pi(\cdot\mid s_t)$ given the past,
\[
  \E_\pi\bigl[\gamma^{t}f(s_t,a_t)\,\Ind[t<T]\bigr]
  =\gamma^{t}\sum_{s}\Pr\nolimits_\pi(s_t=s,\ t<T)\sum_{a}\pi(a\mid s)f(s,a).
\]
Summing over $t\ge0$ and using \eqref{eq:visitation} gives the first equality; the second follows from
$u^\pi=U^\pi d^\pi$. \qedhere
\end{proof}

In a bandit, $u^\pi=d^\pi=U^\pi=1$.

\subsection{The policy gradient theorem and REINFORCE}

Let $\theta\mapsto\pi_\theta$ be a differentiable parametrization of the policies by a vector
$\theta\in\R^n$, and write $\eta(\theta)=\eta(\pi_\theta)$.

\begin{theorem}[Policy gradient theorem \cite{sutton2000,sutton2018}]\label{thm:pg}
With $\pi=\pi_\theta$,
\begin{equation}\label{eq:pg}
  \nabla_\theta\eta(\theta)
  =\sum_{s}u^\pi(s)\sum_{a}\nabla_\theta\pi_\theta(a\mid s)\,Q^\pi(s,a)
  =U^\pi\,\E_{(s,a)\sim\rho^\pi}\bigl[A^\pi(s,a)\,\nabla_\theta\log\pi_\theta(a\mid s)\bigr].
\end{equation}
\end{theorem}

\begin{proof} We refer the reader to the reference provided.
\end{proof}

In particular, the gradient depends on the parametrization only through the score
$\nabla_\theta\log\pi_\theta(a\mid s)$. The REINFORCE algorithm \cite{williams1992} estimates
\eqref{eq:pg} from sampled episodes. With $n_{\mathrm{ep}}$ episodes, visited steps indexed by $m$
(state $s_m$, action $a_m$, time $t_m$, gain $G_m$ from that step) and a baseline $b$, the estimator and
the update are
\begin{equation}\label{eq:reinforce}
  \hat g=\frac1{n_{\mathrm{ep}}}\sum_{m}\gamma^{t_m}\bigl(G_m-b(s_m)\bigr)\nabla_\theta\log\pi_\theta(a_m\mid s_m),
  \qquad
  \theta\leftarrow\theta+\alpha\,\hat g .
\end{equation}
By \cref{lem:visitation} applied to $f(s,a)=\bigl(Q^\pi(s,a)-b(s)\bigr)\nabla\log\pi(a\mid s)$ and
$\E[G_m\mid s_m,a_m]=Q^\pi(s_m,a_m)$, the estimator $\hat g$ is an unbiased estimate of $\nabla_\theta\eta$ whenever
$b$ depends on the state only.

\subsection{The problem: the gradient is covariant, the update is contravariant}\label{sec:problem}

Consider a reparametrization $\theta'=\varphi(\theta)$ of the same policies, with Jacobian
$J=D\varphi(\theta)$, $J_{ij}=\partial\varphi_i/\partial\theta_j$. Writing $\eta(\theta)=\eta'(\varphi(\theta))$, the chain rule gives
\[
  \frac{\partial\eta}{\partial\theta_j}=\sum_i\frac{\partial\eta'}{\partial\theta'_i}\,J_{ij},
\]
and an infinitesimal displacement transforms by linearisation of $\varphi$; thus
\begin{equation}\label{eq:covcontra}
  \nabla_\theta\eta=J\T\,\nabla_{\theta'}\eta,
  \qquad
  \Delta\theta'=J\,\Delta\theta .
\end{equation}
The components of the gradient transform like the basis vectors
$\partial/\partial\theta_j=\sum_iJ_{ij}\,\partial/\partial\theta'_i$: they are \emph{covariant}. The components of a
displacement transform in the opposite way, with $J$ instead of $J\T$: they are \emph{contravariant}.

The plain update $\Delta\theta=\alpha\nabla_\theta\eta$ of \eqref{eq:reinforce} equates a contravariant object with a covariant one,
and this can hold in at most one parametrization. Indeed, apply the same rule in the coordinates $\theta'$:
\[
  \Delta\theta'=\alpha\,\nabla_{\theta'}\eta=\alpha\,J^{-\top}\nabla_\theta\eta .
\]
Read back in the coordinates $\theta$ (to first order),
\begin{equation}\label{eq:plain-in-other-chart}
  \Delta\theta=J^{-1}\Delta\theta'=\alpha\,(J\T J)^{-1}\nabla_\theta\eta,
\end{equation}
which differs from $\alpha\nabla_\theta\eta$ unless $J\T J=I$. For instance, if $\theta'_1=c\,\theta_1$ and the other
coordinates are unchanged, then $J=\diag(c,1,\dots,1)$ and
\begin{equation}\label{eq:c2}
  \Delta\theta_1=\frac{\alpha}{c^{2}}\,\frac{\partial\eta}{\partial\theta_1}:
\end{equation}
the same family of policies, started from the same policy with the same $\alpha$, follows a different
trajectory. Equivalently, the units of the two sides of the update differ, $[\theta_i]$ on the left and
$[\eta]/[\theta_i]$ on the right. Kakade \cite{kakade2001} calls such a rule \emph{noncovariant}.

The parameters merely label the policies, whereas the problem $\max_\pi\eta(\pi)$ concerns the policies themselves.
We therefore regard the set of policies as a manifold and the parametrizations as charts, and look for an update which
is a statement about policies. Converting the covariant gradient into a contravariant displacement requires a
metric on the manifold; we choose one that is intrinsic and measures how much a policy changes.

% =====================================================================
\section{The manifold of policies}\label{sec:manifold}

\subsection{The manifold and the objective}

\begin{definition}[Manifold of policies]\label{def:manifold}
The policy manifold is the set of strictly positive policies, the product over the states of the open simplices
of distributions on $A$,
\[
  \mani=\prod_{s\in S}\Delta(A)^{\circ},
  \qquad
  \Delta(A)^{\circ}=\Bigl\{\pi\in\R^{A}:\ \pi(a)>0\ \ \forall a,\ \ \sum_a\pi(a)=1\Bigr\}.
\]
A point $p\in\mani$ is a policy, written $\pi_p(a\mid s)$. The objective is the function
$\eta:\mani\to\R$, $p\mapsto\eta(\pi_p)$.
\end{definition}

The problem is $\max_{p\in\mani}\eta(p)$; no parametrization is involved. For each state, $p\mapsto\pi_p(\cdot\mid s)$
is a map $\mani\to\Delta(A)^\circ$; Kakade's ``probability manifold of the state $s$'' is the target $\Delta(A)^\circ$ of this map.

\subsection{A redundant parametrization: the softmax}

The standard parametrization of $\mani$ is the softmax with one logit per state--action pair,
$\theta\in\R^{|S||A|}$ with components $\theta_{s,a}$,
\begin{equation}\label{eq:softmax}
  \pi_\theta(a\mid s)=\frac{e^{\theta_{s,a}}}{\sum_{b\in A}e^{\theta_{s,b}}} .
\end{equation}
It is not injective: for $c\in\R^{S}$ and $(\theta+c)_{s,a}=\theta_{s,a}+c_s$,
\begin{equation}\label{eq:softmax-invariance}
  \pi_{\theta+c}(a\mid s)=\frac{e^{\theta_{s,a}}e^{c_s}}{\sum_{b}e^{\theta_{s,b}}e^{c_s}}=\pi_\theta(a\mid s).
\end{equation}
Every policy thus has infinitely many parameters, and the parameter space has more dimensions ($|S||A|$) than the
policies need. Such a parametrization cannot label policies uniquely, and it is not a chart in the sense below.

\subsection{Charts}

\begin{definition}[Chart]\label{def:chart}
A chart of $\mani$ is a diffeomorphism $\Theta:\mani\to U$ onto an open set $U\subset\R^n$. The vector
$\theta=\Theta(p)\in\R^n$, with components $\theta_i$, is the coordinate vector of $p$, and
$\Theta^{-1}:\theta\mapsto\pi_\theta$ is the parametrization. A function $f$ on $\mani$ has the coordinate expression
$\tilde f=f\circ\Theta^{-1}$. A second chart $\Theta'$ defines the transition map and its Jacobian,
\begin{equation}\label{eq:transition}
  \theta'=\varphi(\theta),\quad \varphi=\Theta'\circ\Theta^{-1},\quad
  J=D\varphi(\theta),\quad J_{ij}=\frac{\partial\varphi_i}{\partial\theta_j},\quad
  \tilde f=\tilde f'\circ\varphi.
\end{equation}
\end{definition}

To repair \eqref{eq:softmax}, we fix one reference action $\aref\in A$, set $\Ao=A\setminus\{\aref\}$, and use only the logits
relative to it:
\begin{equation}\label{eq:chart}
  \Theta_{s,a}(p)=\log\pi_p(a\mid s)-\log\pi_p(\aref\mid s),\qquad s\in S,\ a\in\Ao .
\end{equation}

\begin{proposition}\label{prop:chart}
The map $\Theta$ of \eqref{eq:chart} is a chart onto $\R^{|S|(|A|-1)}$, with inverse
\begin{equation}\label{eq:chart-inverse}
  \pi_\theta(a\mid s)=\frac{e^{\theta_{s,a}}}{1+\sum_{b\in\Ao}e^{\theta_{s,b}}}\quad(a\in\Ao),
  \qquad
  \pi_\theta(\aref\mid s)=\frac{1}{1+\sum_{b\in\Ao}e^{\theta_{s,b}}} .
\end{equation}
In particular $\dim\mani=n=|S|(|A|-1)$.
\end{proposition}

\begin{proof}
Both maps are smooth. The right-hand side of \eqref{eq:chart-inverse} is strictly positive and sums to one over $A$, hence
lies in $\mani$. Applying \eqref{eq:chart} to it gives, for $a\in\Ao$,
\[
  \log\frac{e^{\theta_{s,a}}}{1+\sum_{b}e^{\theta_{s,b}}}-\log\frac{1}{1+\sum_{b}e^{\theta_{s,b}}}=\theta_{s,a},
\]
so $\Theta\circ\Theta^{-1}=\mathrm{id}$. Conversely, if $\theta=\Theta(p)$, then $\pi_p(a\mid s)=e^{\theta_{s,a}}\pi_p(\aref\mid s)$ for $a\in\Ao$, and
normalisation $\pi_p(\aref\mid s)\bigl(1+\sum_{b\in\Ao}e^{\theta_{s,b}}\bigr)=1$ determines $\pi_p(\aref\mid s)$; hence
$\Theta^{-1}\circ\Theta=\mathrm{id}$. The image is all of $\R^{|S|(|A|-1)}$ since any real vector $\theta$ is reached by \eqref{eq:chart-inverse}.
\end{proof}

Another reference action $\aref'$ gives another chart. Using $\theta_{s,\aref}:=0$,
\[
  \theta'_{s,a}=\log\pi(a\mid s)-\log\pi(\aref'\mid s)
  =\bigl[\log\pi(a\mid s)-\log\pi(\aref\mid s)\bigr]-\bigl[\log\pi(\aref'\mid s)-\log\pi(\aref\mid s)\bigr],
\]
that is
\begin{equation}\label{eq:change-ref}
  \theta'_{s,a}=\theta_{s,a}-\theta_{s,\aref'},\qquad a\in A\setminus\{\aref'\},
\end{equation}
an affine, non-diagonal transition map. Throughout, $\theta$ without further qualification denotes the coordinates in
the chart \eqref{eq:chart}; the redundant vector of \eqref{eq:softmax} appears only to illustrate its failure.

\subsection{Tangent vectors and covectors}

\begin{definition}[Vectors, covectors, differential]\label{def:tangent}
A tangent vector $v\in T_p\mani$ is the velocity $\dot\gamma(0)$ of a smooth curve $\gamma$ in $\mani$ with $\gamma(0)=p$.
Its components in the chart $\Theta$ are
\[
  v^{\Theta}=\frac{\dd}{\dd t}\Big|_{t=0}\Theta(\gamma(t))\in\R^n ;
\]
the map $v\mapsto v^\Theta$ identifies $T_p\mani$ linearly with $\R^n$, and the basis vector $\partial_i$ is the one with
components $e_i$ (the velocity of the $i$-th coordinate line). A covector is a linear form $\omega$ on $T_p\mani$, with
components $(\omega_\Theta)_i=\omega(\partial_i)$, so that $\omega(v)=\omega_\Theta\T v^\Theta$. The differential of a function
$f$ is the covector
\[
  (\dd f)_p(v)=\frac{\dd}{\dd t}\Big|_{t=0}f(\gamma(t)),
  \qquad
  (\dd f)_p(\partial_i)=\frac{\partial\tilde f}{\partial\theta_i}\bigl(\Theta(p)\bigr).
\]
\end{definition}

By the chain rule, $\frac{\dd}{\dd t}\tilde f(\Theta(\gamma(t)))=D\tilde f\,v^\Theta$, which depends on $\gamma$ only through $v^\Theta$ and is linear
in it; this shows that $\dd f$ is a covector. The usual gradient $\nabla\tilde f$ is thus the component vector of the covector $\dd f$.

\begin{proposition}[Transformation laws]\label{prop:transformation}
Under the change of chart \eqref{eq:transition}, with $J=D\varphi(\Theta(p))$, the components of a vector $v$, of a covector $\omega$ and of a
bilinear form $B$ on $T_p\mani$ transform as
\begin{equation}\label{eq:transformation}
  v^{\Theta'}=J\,v^\Theta,\qquad
  \omega_{\Theta'}=J^{-\top}\omega_\Theta,\qquad
  B'=J^{-\top}B\,J^{-1}.
\end{equation}
\end{proposition}

\begin{proof}
For a curve $\gamma$, $\Theta'(\gamma(t))=\varphi(\Theta(\gamma(t)))$; differentiating at $t=0$ gives $v^{\Theta'}=J\,v^\Theta$.
For a covector, the number $\omega(v)$ does not depend on the chart, so for all $v^\Theta\in\R^n$
\[
  \omega_\Theta\T v^\Theta=\omega_{\Theta'}\T v^{\Theta'}=\omega_{\Theta'}\T J\,v^\Theta
  \quad\Longrightarrow\quad
  \omega_\Theta=J\T\omega_{\Theta'} .
\]
For a bilinear form, $B(v,w)$ does not depend on the chart either, so for all $v^\Theta,w^\Theta$
\[
  (v^\Theta)\T B\,w^\Theta=(Jv^\Theta)\T B'(Jw^\Theta)
  \quad\Longrightarrow\quad
  B=J\T B'J . \qedhere
\]
\end{proof}

\Cref{eq:covcontra} is thereby recovered for intrinsic objects: the displacement is a vector and the differential of $\eta$ is a covector.
The statement $\nabla_\theta\eta=J\T\nabla_{\theta'}\eta$ is the second law of \eqref{eq:transformation} read in the opposite direction.

\begin{definition}[Coordinate derivatives]\label{def:coordder}
Fix a chart $\Theta$. For a function $f$ on $\mani$, let $\partial_if$ and $\partial_i\partial_jf$ be the functions on $\mani$ with
coordinate expressions
\begin{equation}\label{eq:coordder}
  (\partial_if)(q)=(\dd f)_q\bigl(\partial_i|_q\bigr)=\frac{\partial\tilde f}{\partial\theta_i}\bigl(\Theta(q)\bigr),
  \qquad
  (\partial_i\partial_jf)(q)=\frac{\partial^2\tilde f}{\partial\theta_i\,\partial\theta_j}\bigl(\Theta(q)\bigr).
\end{equation}
\end{definition}

In words, $\partial_i$ means ``differentiate the coordinate expression with respect to $\theta_i$ and regard the result as a function on
$\mani$''; the chart is implicit in the symbol. The numbers $(\partial_if)(p)$ are the components of the covector $(\dd f)_p$, whereas the matrix
$(\partial_i\partial_jf)(p)$ is in general not the matrix of a bilinear form, except at a critical point (\cref{prop:hessian}).

% =====================================================================
\section{The Fisher metric and the natural gradient}\label{sec:metric}

\subsection{The Kullback--Leibler divergence}

\begin{definition}[Kullback--Leibler divergence]\label{def:kl}
For two distributions $x,y\in\Delta(A)$,
\begin{equation}\label{eq:kl}
  \KL(x\,\|\,y)=\sum_{a}x_a\log\frac{x_a}{y_a},
\end{equation}
with the conventions $0\log\frac0y=0$ for every $y\ge0$ and $x\log\frac x0=+\infty$ for $x>0$.
\end{definition}

The divergence measures how much $y$ differs from $x$, but it is not a metric in the topological sense: it is not symmetric,
$\KL(x\|y)\ne\KL(y\|x)$ in general, and it does not satisfy the triangle inequality. It does satisfy the following property.

\begin{proposition}[Gibbs' inequality \cite{gibbs}]\label{prop:gibbs}
$\KL(x\|y)\ge0$, with equality if and only if $x=y$.
\end{proposition}

\begin{proof}
From $\log t\le t-1$, with equality iff $t=1$, and writing $a$ for the indices with $x_a>0$,
\[
  -\KL(x\|y)=\sum_{a:\,x_a>0}x_a\log\frac{y_a}{x_a}
  \le\sum_{a:\,x_a>0}(y_a-x_a)=\sum_{a:\,x_a>0}y_a-1\le0 .
\]
Equality requires $y_a=x_a$ whenever $x_a>0$ and $\sum_{x_a>0}y_a=1$, so that $y_a=0=x_a$ elsewhere. \qedhere
\end{proof}

On $\mani$ all probabilities are positive and $\KL$ is finite. Fix a base policy $p\in\mani$ and a state $s$, and define two families of
functions on $\mani$,
\begin{equation}\label{eq:ell-D}
  \ell_{s,a}(q)=\log\pi_q(a\mid s),
  \qquad
  D_s(q)=\KL\bigl(\pi_p(\cdot\mid s)\,\big\|\,\pi_q(\cdot\mid s)\bigr)
  =\sum_a\pi_p(a\mid s)\bigl[\log\pi_p(a\mid s)-\ell_{s,a}(q)\bigr].
\end{equation}
By \cref{prop:gibbs}, $D_s\ge0$, with $D_s(q)=0$ iff $\pi_q(\cdot\mid s)=\pi_p(\cdot\mid s)$. Thus $D_s$ has a global minimum at $q=p$, and
\begin{equation}\label{eq:D-critical}
  D_s(p)=0,\qquad(\dd D_s)_p=0 .
\end{equation}
The expansion of $D_s$ around $p$ therefore starts at second order, which suggests using the second derivative of $D_s$ as a measure of the change
of the policy. For this the second derivative must be an intrinsic object.

\subsection{The Hessian at a critical point}

\begin{proposition}[Intrinsic Hessian]\label{prop:hessian}
Let $f$ be a function on $\mani$ with $(\dd f)_p=0$, and let $H_{ij}=(\partial_i\partial_jf)(p)$ in a chart $\Theta$. Under a change of chart \eqref{eq:transition},
\begin{equation}\label{eq:hessian-law}
  H'=J^{-\top}H\,J^{-1},
\end{equation}
so that $\Hess_pf(v,w):=\sum_{i,j}H_{ij}v^\Theta_iw^\Theta_j$ is a well-defined symmetric bilinear form on $T_p\mani$, independent of the chart.
\end{proposition}

\begin{proof}
Differentiating $\tilde f=\tilde f'\circ\varphi$ once,
\[
  \frac{\partial\tilde f}{\partial\theta_i}=\sum_k\frac{\partial\tilde f'}{\partial\theta'_k}\bigl(\varphi(\theta)\bigr)\frac{\partial\varphi_k}{\partial\theta_i},
\]
and once more, by the product and chain rules,
\[
  \frac{\partial^2\tilde f}{\partial\theta_i\partial\theta_j}
  =\sum_{k,l}\frac{\partial^2\tilde f'}{\partial\theta'_k\partial\theta'_l}\,J_{kj}J_{li}\;+\;\sum_k\frac{\partial\tilde f'}{\partial\theta'_k}\,\frac{\partial^2\varphi_k}{\partial\theta_i\partial\theta_j}.
\]
In matrix form,
\begin{equation}\label{eq:hessian-chain}
  D^2\tilde f(v,w)=D^2\tilde f'(Jv,Jw)+D\tilde f'\bigl(D^2\varphi(v,w)\bigr).
\end{equation}
The second term is proportional to the first derivative of $f$ and vanishes at a critical point. What remains is $H=J\T H'J$, which is \eqref{eq:hessian-law}, the
transformation law of a bilinear form in \eqref{eq:transformation}.
\end{proof}


\subsection{The Fisher information of a state}

\begin{definition}[Fisher information of a state]\label{def:fisher-state}
Given the base policy $p$ and the state $s$, the Fisher information is the Hessian of the Kullback--Leibler divergence at its minimum,
\[
  F_s(v,w)=\Hess_pD_s(v,w),\qquad F_s(\partial_i,\partial_j)=(\partial_i\partial_jD_s)(p).
\]
\end{definition}

By \eqref{eq:D-critical} and \cref{prop:hessian}, $F_s$ is a chart-independent bilinear form, and in any chart
\begin{equation}\label{eq:kl-quadratic}
  \KL\bigl(\pi_p(\cdot\mid s)\,\big\|\,\pi_{\Theta^{-1}(\theta+\delta)}(\cdot\mid s)\bigr)=\tfrac12F_s(\delta,\delta)+O(\|\delta\|^3):
\end{equation}
steps measured by $F_s$ are steps measured in Kullback--Leibler divergence.

\begin{proposition}[Computation of $F_s$]\label{prop:fisher-score}
\begin{equation}\label{eq:fisher-score}
  F_s(\partial_i,\partial_j)=\E_{a\sim\pi_p(\cdot\mid s)}\bigl[(\partial_i\ell_{s,a})(p)\,(\partial_j\ell_{s,a})(p)\bigr].
\end{equation}
\end{proposition}

\begin{proof}
By \eqref{eq:ell-D}, the weights $\pi_p(a\mid s)$ being fixed numbers,
\[
  \partial_i\partial_jD_s=-\sum_a\pi_p(a\mid s)\,\partial_i\partial_j\ell_{s,a}.
\]
Since $\partial_i\partial_j\log\pi=\pi^{-1}\partial_i\partial_j\pi-\partial_i\log\pi\,\partial_j\log\pi$, evaluation at $p$ gives
\[
  \partial_i\partial_jD_s(p)=-\sum_a\partial_i\partial_j\pi_p(a\mid s)+\sum_a\pi_p(a\mid s)\,\partial_i\ell_{s,a}(p)\,\partial_j\ell_{s,a}(p).
\]
The first sum vanishes, as $\sum_a\pi_\theta(a\mid s)=1$ for all $\theta$ and hence $\sum_a\partial_i\partial_j\pi_\theta(a\mid s)=0$.
\end{proof}

In particular $F_s(v,v)=\E_a\bigl[((\dd\ell_{s,a})_p(v))^2\bigr]\ge0$: the Fisher information is positive semidefinite, and the covector $(\dd\ell_{s,a})_p$ is the score.
The formula is manifestly of the form ``average of a product of covectors'', which gives a second proof of the transformation law $F_s'=J^{-\top}F_sJ^{-1}$.

\paragraph{The Fisher matrix in the chart \eqref{eq:chart}.}
Put $z_{s,a}=\theta_{s,a}$ for $a\in\Ao$ and $z_{s,\aref}=0$. Then
\[
  \ell_{s,a}=z_{s,a}-\log\sum_ce^{z_{s,c}},
\]
and for $b\in\Ao$,
\[
  \frac{\partial z_{s,a}}{\partial\theta_{s',b}}=\delta_{ss'}\Ind[a=b],
  \qquad
  \frac{\partial}{\partial\theta_{s',b}}\log\sum_ce^{z_{s,c}}=\delta_{ss'}\frac{e^{z_{s,b}}}{\sum_ce^{z_{s,c}}}=\delta_{ss'}\,\pi(b\mid s),
\]
where $\Ind[a=b]=0$ if $a=\aref$. Hence
\begin{equation}\label{eq:score}
  \partial_{s',b}\,\ell_{s,a}=\delta_{ss'}\bigl(\Ind[a=b]-\pi(b\mid s)\bigr),\qquad b\in\Ao .
\end{equation}
Insert \eqref{eq:score} into \eqref{eq:fisher-score}; write $\pi_a=\pi_p(a\mid s)$ and take $b,c\in\Ao$. Then
\begin{align*}
  F_s(\partial_{s,b},\partial_{s,c})
  &=\sum_a\pi_a\bigl(\Ind[a=b]-\pi_b\bigr)\bigl(\Ind[a=c]-\pi_c\bigr)\\
  &=\pi_b\delta_{bc}-\pi_c\pi_b-\pi_b\pi_c+\pi_b\pi_c\sum_a\pi_a\\
  &=\pi_b\delta_{bc}-\pi_b\pi_c,
\end{align*}
and $F_s(\partial_{s,b},\partial_{s',c})=0$ for $s'\ne s$, because the score of $s$ does not depend on the logits of the other states.

\begin{proposition}[The Fisher matrix of a state is positive definite]\label{prop:fisher-pd}
In the chart \eqref{eq:chart}, with $\pio=(\pi_p(b\mid s))_{b\in\Ao}$,
\begin{equation}\label{eq:Cs}
  F_s=C_s,\qquad C_s=\diag(\pio)-\pio\pio\T,\qquad
  v\T C_sv=\Var_{a\sim\pi_p(\cdot\mid s)}\bigl[v_a\bigr]\quad(v_{\aref}:=0),
\end{equation}
and $C_s$ is positive definite.
\end{proposition}

\begin{proof}
By \eqref{eq:score}, $(\dd\ell_{s,a})(v)=v_a-\sum_{b\in\Ao}\pi_bv_b=v_a-\E_{a'\sim\pi}[v_{a'}]$, using $v_{\aref}=0$, and the average of its square over $a\sim\pi$
is the variance. If the variance vanishes, $a\mapsto v_a$ is constant on $A$ (all $\pi_a>0$), and the constant is $v_{\aref}=0$; hence $v=0$.
\end{proof}

\paragraph{A counterexample: the redundant softmax.}
Compute the same objects in the parameter $\theta\in\R^{|S||A|}$ of \eqref{eq:softmax}, which is not a chart. Now
$\ell_{s,a}(\theta)=\theta_{s,a}-\log\sum_be^{\theta_{s,b}}$ depends only on the block $\theta_s$, and
\[
  \frac{\partial\ell_{s,a}}{\partial\theta_{s,b}}=\Ind[a=b]-\pi(b\mid s)\quad\text{for all }b\in A,
  \qquad
  F_s^{\mathrm{red}}=\E_a\bigl[\nabla\ell\,\nabla\ell\T\bigr]=\diag(\pi_s)-\pi_s\pi_s\T,
\]
with $\pi_s=(\pi_p(b\mid s))_{b\in A}$, by the same computation. Let $v=\mathbf 1_s$ be the vector with entries $1$ on the block of $s$ and $0$ elsewhere. For every $a$,
\[
  (\dd\ell_{s,a})(v)=\sum_b\frac{\partial\ell_{s,a}}{\partial\theta_{s,b}}=\sum_b\bigl(\Ind[a=b]-\pi(b\mid s)\bigr)=1-1=0,
\]
so $F^{\mathrm{red}}_s(v,v)=0$ although $v\ne0$; equivalently $F^{\mathrm{red}}_s\mathbf 1=\pi_s-\pi_s=0$. This is the invariance \eqref{eq:softmax-invariance}
at the infinitesimal level: $\ell_{s,a}(\theta+t\mathbf 1_s)=\ell_{s,a}(\theta)$ for all $t$. The matrix $F^{\mathrm{red}}_s$ is singular and assigns zero
length to a nonzero displacement, so it is not a metric. The cause is the non-injectivity of the parametrization: the differential of $\theta\mapsto\pi_\theta$ has a kernel, and
the Fisher matrix is the pullback of a positive form on the policies. In the chart \eqref{eq:chart} no such direction exists, since the coordinate $\theta_{s,\aref}$ is pinned. A bijective
chart is thus necessary for a nondegenerate metric.

\begin{remark}[Necessary but not sufficient: damping]\label{rem:damping}
Bijectivity removes the kernel that comes from the labelling. Positive definiteness of the Fisher metric of the whole policy (\eqref{eq:fisher-policy} below) requires in
addition that no direction of $T_p\mani$ leaves $\pi_p(\cdot\mid s)$ unchanged at all states of positive visitation. In the tabular chart this holds iff every state has
$d_p(s)>0$, but for restricted families of policies (linear features, neural networks) it is an identifiability condition which is hard to verify. Even when it holds, the matrix can be
numerically singular, since the weight $d_p(s)$ of rarely visited states is small, and a matrix estimated from finitely many samples is singular when some states or actions are not
visited. In practice one therefore damps the metric: $F$ is replaced by $F+\lambda\Lambda$ with $\lambda>0$ and $\Lambda$ positive definite (Tikhonov regularisation, as in the
Levenberg--Marquardt method). This is a numerical device: it is not part of the geometry, it depends on the chart through $\Lambda$, and it changes the metric.
\Cref{sec:regularisation} shows that the way it is applied matters.
\end{remark}

\subsection{The Fisher metric of the policy}

\begin{definition}[Fisher metric of the policy]\label{def:fisher-policy}
The Fisher metric of the policy $p$ is the average of the metrics of the states over the discounted visitation $d_p$ of \eqref{eq:visitation}, held fixed (not differentiated),
\begin{equation}\label{eq:fisher-policy}
  F_p=\sum_sd_p(s)\,F_s,
  \qquad
  F_p(v,w)=\E_{(s,a)\sim\rho_p}\bigl[(\dd\ell_{s,a})_p(v)\,(\dd\ell_{s,a})_p(w)\bigr].
\end{equation}
\end{definition}

It is a sum of forms on the same space $T_p\mani$, hence an intrinsic, positive semidefinite bilinear form. By \cref{prop:fisher-pd} it is positive definite in the chart
\eqref{eq:chart} as soon as all states have $d_p(s)>0$, and it is block diagonal with blocks $d_p(s)C_s$. Positive definite bilinear forms are metrics, and a metric allows one to convert a
covector into a vector.

\begin{definition}[Gradient with respect to a metric]\label{def:grad}
For a metric $g$ with matrix $[g]_{ij}=g(\partial_i,\partial_j)$, which transforms by \eqref{eq:transformation}, the gradient of $f$ is the vector $\grad_gf\in T_p\mani$ defined by
\begin{equation}\label{eq:grad}
  g(\grad_gf,v)=(\dd f)_p(v)\ \ \forall v\in T_p\mani,
  \qquad\text{i.e.}\qquad
  (\grad_gf)^\Theta=[g]^{-1}\nabla\tilde f .
\end{equation}
\end{definition}

\begin{proposition}[Steepest ascent]\label{prop:steepest}
Among the displacements of fixed $g$-length $g(v,v)=\delta^2$, the function $v\mapsto(\dd f)_p(v)$ is maximised by
\[
  v=\delta\,\frac{\grad_gf}{\|\grad_gf\|_g},
\]
with maximal value $\delta\|\grad_gf\|_g=\delta\sqrt{\nabla\tilde f\T[g]^{-1}\nabla\tilde f}$. In components the direction is proportional to $[g]^{-1}\nabla\tilde f$, so it depends on the metric.
\end{proposition}

\begin{proof}
By \eqref{eq:grad}, $(\dd f)_p(v)=g(\grad_gf,v)$. Cauchy--Schwarz for the inner product $g$ gives
\[
  g(\grad_gf,v)\le\|\grad_gf\|_g\,\|v\|_g=\delta\,\|\grad_gf\|_g,
\]
with equality iff $v$ is a positive multiple of $\grad_gf$. The value of the norm follows from $\|\grad_gf\|_g^2=\nabla\tilde f\T[g]^{-1}[g][g]^{-1}\nabla\tilde f$.
\end{proof}

\paragraph{Covariance.}
The vector $\grad_gf$ is defined without reference to a chart, so its components transform with $J$. This can be verified directly from \eqref{eq:transformation}: with $[g']=J^{-\top}[g]J^{-1}$ and $\nabla\tilde f'=J^{-\top}\nabla\tilde f$,
\begin{equation}\label{eq:covariance}
  [g']^{-1}\nabla\tilde f'=\bigl(J[g]^{-1}J\T\bigr)\bigl(J^{-\top}\nabla\tilde f\bigr)=J\,[g]^{-1}\nabla\tilde f .
\end{equation}
Plain gradient ascent is the choice $[g]=I$ in the chart in use. The matrix of that metric in another chart is $J^{-\top}J^{-1}\ne I$, so the algorithm is a rule ``chart $\mapsto$ metric'' and
changes with the chart; this is \eqref{eq:plain-in-other-chart}. A covariant algorithm needs a single metric on $\mani$, and the Fisher metric is one which, moreover, measures
Kullback--Leibler divergence.

\begin{definition}[Natural direction and KL budget]\label{def:natural}
The natural direction is $v_\star=\bigl(\dd\eta/U_p\bigr)^\sharp$ with respect to $F_p$, the factor $U_p$ being the normalisation between the weights $u_p$ of $\dd\eta$ in \eqref{eq:pg} and
$d_p$ of $F_p$; only the direction matters. Its components, the weight vector $w=v_\star^\Theta\in\R^n$, solve
\begin{equation}\label{eq:natural-eq}
  F\,w=\bar g,\qquad\bar g=\frac{\nabla\tilde\eta}{U_p}.
\end{equation}
For a prescribed Kullback--Leibler budget $\varepsilon$ per update, the update is
\begin{equation}\label{eq:budget}
  \theta\leftarrow\theta+\alpha\,w,\qquad\alpha=\sqrt{\frac{2\varepsilon}{w\T Fw}} .
\end{equation}
\end{definition}

By \eqref{eq:kl-quadratic} and \eqref{eq:fisher-policy}, the update changes the policy by a Kullback--Leibler divergence of about $\tfrac12\alpha^2w\T Fw=\varepsilon$, averaged over the states with weights
$d_p$. The number $\varepsilon$ is dimensionless and means the same in every chart, and by \eqref{eq:covariance} the direction is covariant.

\paragraph{Flow versus discrete step.}
The continuous-time flow $\dot\gamma=\grad_g\eta(\gamma)$ is an equation between vectors of $T_{\gamma(t)}\mani$ and is exactly chart-independent. The discrete step $\theta\mapsto\theta+\alpha w$, in contrast, stands for
\[
  \mathrm{Step}_\Theta(p,v)=\Theta^{-1}\bigl(\Theta(p)+\alpha v^\Theta\bigr),
\]
since $\mani$ itself has no addition. Using another chart $\Theta'$ with transition map $\varphi$ gives
\[
  \Theta\bigl(\mathrm{Step}_{\Theta'}(p,v)\bigr)=\varphi^{-1}\bigl(\varphi(\theta)+\alpha J v^\Theta\bigr)=\theta+\alpha v^\Theta+O(\alpha^2),
\]
so the two steps agree to first order and can differ at $O(\alpha^2)$ unless $\varphi$ is affine.

The discrete natural gradient is therefore covariant to first order in $\alpha$ only; the intrinsic alternative is to follow a curve in $\mani$, for instance a geodesic of $F_p$.

\subsection{The natural gradient in closed form}

In the chart \eqref{eq:chart}, \eqref{eq:pg} and \eqref{eq:score} give the components of $\dd\eta$. For $b\in\Ao$,
\begin{align}
  (\dd\eta)_p(\partial_{s,b})
  &=u_p(s)\sum_a\pi(a\mid s)\,Q(s,a)\bigl(\Ind[a=b]-\pi(b\mid s)\bigr)\nonumber\\
  &=u_p(s)\Bigl[\pi(b\mid s)Q(s,b)-\pi(b\mid s)\sum_a\pi(a\mid s)Q(s,a)\Bigr]\nonumber\\
  &=u_p(s)\,\pi(b\mid s)\,A(s,b).\label{eq:grad-closed}
\end{align}

\begin{proposition}[Natural gradient in the chart]\label{prop:natural-closed}
The weight vector of \eqref{eq:natural-eq} is
\begin{equation}\label{eq:natural-closed}
  w_{s,a}=A^\pi(s,a)-A^\pi(s,\aref),\qquad a\in\Ao,
\end{equation}
and
\[
  w\T Fw=\sum_sd_p(s)\,\Var_{a\sim\pi(\cdot\mid s)}\bigl[A^\pi(s,a)\bigr].
\]
\end{proposition}

\begin{proof}
By \eqref{eq:Cs} and \eqref{eq:grad-closed}, the block of $s$ of $Fw=\bar g$ reads
\[
  d_p(s)\,C_sw_s=d_p(s)\,\pio\odot A_\circ,
\]
that is, componentwise for $b\in\Ao$,
\[
  \pi_b\bigl(w_b-\pio\cdot w_\circ\bigr)=\pi_bA_b .
\]
Hence $w_b=A_b+c_s$ with $c_s=\pio\cdot w_\circ$. Taking the scalar product of $w_\circ=A_\circ+c_s\mathbf 1$ with $\pio$ and using $\sum_a\pi_aA_a=0$,
\[
  c_s=\sum_{b\in\Ao}\pi_bA_b+c_s(1-\pi_{\aref})=-\pi_{\aref}A_{\aref}+c_s(1-\pi_{\aref})
  \quad\Longrightarrow\quad c_s=-A_{\aref}.
\]
Finally $w\T Fw=\sum_sd_p(s)\Var_\pi[w_s]$ by \eqref{eq:Cs}, and a constant shift does not change the variance, so $\Var_\pi[w_s]=\Var_\pi[A(s,\cdot)]$.
\end{proof}

The gradient \eqref{eq:grad-closed} carries the factors $u_p(s)$ and $\pi(b\mid s)$: it is small for rarely visited states and for nearly deterministic policies. The natural direction
\eqref{eq:natural-closed} has neither factor, the weights $d_p(s)$ and $\pi$ being cancelled by $F^{-1}$.

\paragraph{The resulting policy update.}
For a step of multiplier $\alpha$ the log-odds of $a$ relative to $\aref$ change by $\alpha\bigl(A(s,a)-A(s,\aref)\bigr)$. Hence
\[
  \frac{\pi'(a\mid s)}{\pi'(\aref\mid s)}=\frac{\pi(a\mid s)}{\pi(\aref\mid s)}\,e^{\alpha(A(s,a)-A(s,\aref))}
  =\frac{\pi(a\mid s)\,e^{\alpha A(s,a)}}{\pi(\aref\mid s)\,e^{\alpha A(s,\aref)}},
\]
and normalisation gives
\begin{equation}\label{eq:exp-update}
  \pi'(a\mid s)=\frac{\pi(a\mid s)\,e^{\alpha A^\pi(s,a)}}{\sum_b\pi(b\mid s)\,e^{\alpha A^\pi(s,b)}},
\end{equation}
an exponentiated-advantage update, independent of the reference action, whose limit $\alpha\to\infty$ is greedy policy improvement.

In a bandit the formulas read
\[
  \frac{\partial\eta}{\partial\theta_b}=\pi_b(\mu_b-\eta),\qquad
  w_b=\mu_b-\mu_{\aref},\qquad
  w\T Fw=\Var_{a\sim\pi}[\mu_a].
\]

% =====================================================================
\section{The algorithm}\label{sec:algorithm}

\subsection{Features and Kakade's theorem}

\begin{definition}[Features and compatible approximator]\label{def:features}
The feature vector of a state--action pair is the score in the chart,
$\xi(s,a)=\bigl(\partial_i\ell_{s,a}\bigr)_i=\nabla_\theta\log\pi_\theta(a\mid s)\in\R^n$. By \eqref{eq:score}, in the chart \eqref{eq:chart},
\begin{equation}\label{eq:features}
  \xi(s,a)=e_s\otimes\chi(s,a),\qquad\chi(s,a)=\iota(a)-\pio(s)\in\R^{\Ao},
\end{equation}
where $\iota(a)$ is the indicator vector of $a$ (zero if $a=\aref$), $\pio(s)=(\pi(b\mid s))_{b\in\Ao}$, and $e_s\otimes(\cdot)$ places a vector in the block of the state $s$. For a weight vector $w\in\R^n$
the compatible approximator is the function
\begin{equation}\label{eq:compatible}
  f^w(s,a)=w\T\xi(s,a)=(\dd\ell_{s,a})(v),\qquad v^\Theta=w .
\end{equation}
\end{definition}

By \eqref{eq:fisher-policy} and \eqref{eq:pg},
\[
  F=\E_\rho\bigl[\xi\xi\T\bigr],\qquad\bar g=\E_\rho\bigl[A^\pi\,\xi\bigr],\qquad\rho=\rho_p.
\]

\begin{theorem}[Kakade \cite{kakade2001}]\label{thm:kakade}
Let $\epsilon(w)=\E_{(s,a)\sim\rho}\bigl[\bigl(f^w(s,a)-A^\pi(s,a)\bigr)^2\bigr]$. Then $w$ minimises $\epsilon$ if and only if $Fw=\bar g$. In particular, when $F$ is positive definite, the unique minimiser is
the weight vector of the natural direction, $w=F^{-1}\bar g$.
\end{theorem}

\begin{proof}
Expanding the square,
\[
  \epsilon(w)=w\T\E_\rho[\xi\xi\T]\,w-2\,w\T\E_\rho[A^\pi\xi]+\E_\rho\bigl[(A^\pi)^2\bigr]
  =w\T Fw-2\,w\T\bar g+\E_\rho\bigl[(A^\pi)^2\bigr],
\]
a convex quadratic whose gradient $2(Fw-\bar g)$ vanishes exactly at the solutions of \eqref{eq:natural-eq}.
\end{proof}

Thus the natural gradient is the solution of a least-squares regression of the advantage on the features, and no matrix has to be inverted explicitly. Two remarks follow.

First, $\E_\rho[b(s)\xi]=0$ for every function $b$ of the state alone, because the score has zero mean; replacing the target $A^\pi$ by $Q^\pi$ or by $Q^\pi-b(s)$ therefore does not change the minimiser, which is why a baseline is allowed.

Second, in the tabular chart the regression is exact. For $w_{s,a}=A(s,a)-A(s,\aref)$ with $a\in\Ao$ (and $w_{s,\aref}:=0$), \eqref{eq:features} gives
\begin{align*}
  f^w(s,a)&=w_a-\sum_{b\in\Ao}\pi_bw_b,\\
  \sum_{b\in\Ao}\pi_bw_b&=\sum_{b\in\Ao}\pi_bA_b-A_{\aref}(1-\pi_{\aref})
  =-\pi_{\aref}A_{\aref}-A_{\aref}+\pi_{\aref}A_{\aref}=-A_{\aref},
\end{align*}
so that $f^w(s,a)=(A_a-A_{\aref})+A_{\aref}=A(s,a)$ for $a\in\Ao$, and $f^w(s,\aref)=0+A_{\aref}=A(s,\aref)$ as well. Hence $\epsilon(w)=0$.

\subsection{Estimators from samples}

Sample $n_{\mathrm{ep}}$ episodes with the current policy and index the visited steps by $m$: state $s_m$, action $a_m$, time $t_m$ and gain $G_m$ from that step. Set $\hat A_m=G_m-b(s_m)$ for a baseline $b$ depending on
the state only, $\xi_m=\xi(s_m,a_m)$ and $\chi_m=\chi(s_m,a_m)$. Since $\nabla_\theta\log\pi=\xi$, the REINFORCE estimator \eqref{eq:reinforce} of the components of $\dd\eta$ is
\begin{equation}\label{eq:ghat}
  \hat g=\frac1{n_{\mathrm{ep}}}\sum_m\gamma^{t_m}\hat A_m\,\xi_m .
\end{equation}
The matrix $F$ is block diagonal over the states with blocks $d_p(s)C_s$ by \eqref{eq:Cs}, and $\bar g$ has blocks $d_p(s)\,b_s$, where by \eqref{eq:grad-closed}
\begin{equation}\label{eq:Cs-bs}
  C_s=\E_{a\sim\pi(\cdot\mid s)}\bigl[\chi\chi\T\bigr]=\diag(\pio)-\pio\pio\T,
  \qquad
  b_s=\E_{a\sim\pi(\cdot\mid s)}\bigl[A\,\chi\bigr]=\pio\odot A_\circ .
\end{equation}
The equation $Fw=\bar g$ therefore decouples by state, and the visitation cancels:
\begin{equation}\label{eq:decoupled}
  d_p(s)\,C_sw_s=d_p(s)\,b_s\iff C_sw_s=b_s .
\end{equation}
Let $M_s=\{m:s_m=s\}$ be the visits to $s$. The within-state weights, the visitation estimate and the local estimates of $C_s$ and $b_s$ are
\begin{align}
  \omega^{(s)}_m&=\frac{\gamma^{t_m}}{\sum_{m'\in M_s}\gamma^{t_{m'}}}, &
  \hat d_s&=\frac{\sum_{m\in M_s}\gamma^{t_m}}{\sum_m\gamma^{t_m}},\nonumber\\
  \hat C_s&=\sum_{m\in M_s}\omega^{(s)}_m\chi_m\chi_m\T, &
  \hat b_s&=\sum_{m\in M_s}\omega^{(s)}_m\hat A_m\chi_m .\label{eq:local-estimates}
\end{align}
In a bandit there is one state, $\hat d_s=1$ and all weights are equal.

\subsection{Regularisation: why it must be statewise}\label{sec:regularisation}

The local systems $\hat C_s w_s=\hat b_s$ are singular or very noisy when a state has few visits, so they must be regularised (\cref{rem:damping}). The straightforward choice is a ridge on the global system,
$(\hat F+\lambda\Lambda)w=\hat{\bar g}$, whose blocks are $\hat d_s\hat C_s$ and $\hat d_s\hat b_s$. This breaks the cancellation \eqref{eq:decoupled}.

\begin{remark}[A global ridge destroys the cancellation of the visitation]\label{rem:global-ridge}
For the block of the state $s$, with $\Lambda_s$ the corresponding block of $\Lambda$,
\[
  \bigl(\hat d_s\hat C_s+\lambda\Lambda_s\bigr)w_s=\hat d_s\hat b_s
  \iff
  \Bigl(\hat C_s+\frac{\lambda}{\hat d_s}\Lambda_s\Bigr)w_s=\hat b_s .
\]
The effective ridge $\lambda/\hat d_s$ grows without bound as $\hat d_s\to0$. For a rarely visited state it dominates $\hat C_s$, and
\[
  w_s\approx\frac{\hat d_s}{\lambda}\,\Lambda_s^{-1}\hat b_s,
\]
which is proportional to the visitation $\hat d_s$: exactly the factor that the natural gradient is designed to cancel (\cref{prop:natural-closed}), and the behaviour of plain gradient ascent.
The effect does not disappear when the number of samples grows, since $\hat d_s$ is a property of the state, not of the sample size.
\end{remark}

The ridge must therefore be applied to the conditional system \eqref{eq:decoupled}, state by state, with a strength that depends on the amount of data of that state and not on its visitation.
With the Kish effective sample size of the visits to $s$,
\begin{equation}\label{eq:neff}
  n_{\mathrm{eff}}(s)=\frac{\bigl(\sum_{m\in M_s}\gamma^{t_m}\bigr)^2}{\sum_{m\in M_s}\gamma^{2t_m}},
\end{equation}
a ridge constant $\lambda_0>0$ and a shrinkage constant $n_0\ge0$, the statewise weight vector is
\begin{equation}\label{eq:statewise}
  \lambda_s=\frac{\lambda_0}{\sqrt{n_{\mathrm{eff}}(s)}},\qquad
  c_s=\frac{n_{\mathrm{eff}}(s)}{n_{\mathrm{eff}}(s)+n_0},\qquad
  \hat w_s=c_s\bigl(\hat C_s+\lambda_s\Lambda_s\bigr)^{-1}\hat b_s,
\end{equation}
with $\Lambda_s=I$ in the chart $\Theta$ (see \eqref{eq:chart-psi} for other charts) and $\hat w_s=0$ for an unvisited state. As the data of a state accumulate, $\lambda_s\to0$ and $c_s\to1$, so
$\hat w_s\to C_s^{-1}b_s=w_s$ of \eqref{eq:natural-closed}: the regularisation vanishes where the data are plentiful, and it shrinks the update where they are not. The ridge is a statistical device and is
not part of the metric, so it does not enter the Kullback--Leibler norm:
\begin{equation}\label{eq:q-alpha}
  q=\sum_s\hat d_s\,\hat w_s\T\hat C_s\hat w_s\approx w\T Fw,\qquad
  \alpha=\sqrt{\frac{2\varepsilon}{q}},\qquad
  \theta\leftarrow\theta+\alpha\,\hat w,
\end{equation}
with $\alpha=0$ if $q$ is numerically zero. Plain gradient ascent is $\theta\leftarrow\theta+\mathrm{lr}\cdot\hat g$ with \eqref{eq:ghat}.

\subsection{Step size and safeguards}

\paragraph{Horizon bias.}
If an episode is cut after $T_{\max}$ steps, the gain of its last step is bootstrapped with the exact value, $G=V^\pi(s_{T_{\max}})$, so that the estimates target the infinite-horizon objective.

\paragraph{Step cap.}
Near a deterministic policy the variances in \cref{prop:natural-closed} tend to zero, so $q\to0$ and $\alpha\to\infty$. For both methods the update $\Delta\theta$ is rescaled, without changing its direction, so that
\begin{equation}\label{eq:cap}
  \Delta\theta\leftarrow\Delta\theta\cdot\min\Bigl\{1,\ \frac{c_{\max}}{\|\Delta\theta\|_\infty}\Bigr\}.
\end{equation}
The coordinates themselves are not clipped: with the reference logit fixed, clipping could force two nearly equal best actions onto the same bound.

\paragraph{Uninformative batches.}
If a batch carries no information (every pull on one arm in a bandit), then $\hat g=0$ up to round-off and $\alpha=\sqrt{2\varepsilon/q}$ would inflate round-off noise to a step of full size in a random direction.
The update is skipped when $q\le\tau\cdot\frac1N\sum_m\hat A_m^2$ (bandit) or when $q$ is numerically zero (gridworld).

\subsection{The chart in the implementation}\label{sec:chart-impl}

To test covariance, the implementation works in the chart $\Psi=D\Theta$ with coordinates $\psi=D\theta$, where $D=\diag(\delta_i)$, $\delta_i>0$ (the same policies in different units). By \eqref{eq:transformation} with $J=D$,
\begin{equation}\label{eq:chart-psi}
  \xi_\psi=D^{-1}\xi_\theta,\qquad
  \hat C_{s,\psi}=D_s^{-1}\hat C_{s,\theta}D_s^{-1},\qquad
  \hat b_{s,\psi}=D_s^{-1}\hat b_{s,\theta},\qquad
  \Lambda_s=D_s^{-2},
\end{equation}
where $D_s$ is the block of $D$ of the state $s$ and $\Lambda_s=D_s^{-2}$ is the ridge ``$\lambda I$ in the chart $\Theta$'' transported to $\Psi$.

\begin{proposition}[Covariance of the implemented update]\label{prop:covariance-impl}
Take $\Lambda_s=D_s^{-2}$, and measure the cap \eqref{eq:cap} in $\theta$-units, i.e.\ with $\|D^{-1}\Delta\psi\|_\infty\le c_{\max}$. Then $\hat w_{s,\psi}=D_s\hat w_{s,\theta}$, the quantities $q$ and $\alpha$ are the same in both charts, and
$\Delta\theta=D^{-1}\Delta\psi$ is independent of $D$. Plain gradient ascent gives $\Delta\theta=\mathrm{lr}\,D^{-2}\hat g_\theta$, which is \eqref{eq:c2} with $c=\delta_i$.
\end{proposition}

\begin{proof}
By \eqref{eq:chart-psi},
\[
  \bigl(\hat C_{s,\psi}+\lambda_sD_s^{-2}\bigr)=D_s^{-1}\bigl(\hat C_{s,\theta}+\lambda_sI\bigr)D_s^{-1},
\]
hence
\[
  \hat w_{s,\psi}=c_s\,D_s\bigl(\hat C_{s,\theta}+\lambda_sI\bigr)^{-1}D_s\,D_s^{-1}\hat b_{s,\theta}=D_s\,\hat w_{s,\theta},
\]
where $n_{\mathrm{eff}}$, $\lambda_s$ and $c_s$ do not depend on the chart. Next,
\[
  \hat w_{s,\psi}\T\hat C_{s,\psi}\hat w_{s,\psi}=\hat w_{s,\theta}\T D_s\,D_s^{-1}\hat C_{s,\theta}D_s^{-1}\,D_s\hat w_{s,\theta}=\hat w_{s,\theta}\T\hat C_{s,\theta}\hat w_{s,\theta},
\]
so $q$ and $\alpha$ coincide, $\Delta\psi=\alpha\hat w_\psi=D\,\alpha\hat w_\theta$ and $\Delta\theta=D^{-1}\Delta\psi=\alpha\hat w_\theta$. The cap acts on $D^{-1}\Delta\psi=\Delta\theta$ and is therefore the same.
For plain gradient ascent, $\hat g_\psi=D^{-1}\hat g_\theta$ by \eqref{eq:chart-psi}, so $\Delta\psi=\mathrm{lr}\,D^{-1}\hat g_\theta$ and $\Delta\theta=D^{-1}\Delta\psi=\mathrm{lr}\,D^{-2}\hat g_\theta$.
\end{proof}

A change of reference action is a non-diagonal linear change of chart (\eqref{eq:change-ref}): the exact natural direction is a vector of $T_p\mani$ and does not change, whereas the plain gradient step does. The cap
\eqref{eq:cap} and a ridge $\lambda I$ refer to the coordinates in use and are not invariant under such a change.

\subsection{The algorithm and concluding remarks}

\begin{algorithm}[t]
\caption{One sampled update (plain or natural).}\label{alg:update}
\begin{algorithmic}[1]
\Require chart coordinates $\psi$, scaling $D$; method with $\mathrm{lr}$ or $\varepsilon$; $\lambda_0,n_0,c_{\max},n_{\mathrm{ep}},T_{\max}$
\State $\theta\gets D^{-1}\psi$; policy $\pi_\theta$ by \eqref{eq:chart-inverse}; exact $V^\pi$ (baseline and bootstrap)
\State sample $n_{\mathrm{ep}}$ episodes (cut at $T_{\max}$, bootstrapped); gains $G_m$
\State $\hat A_m\gets G_m-V^\pi(s_m)$;\quad $\chi_m\gets D_{s_m}^{-1}\bigl(\iota(a_m)-\pio(s_m)\bigr)$ \Comment{\eqref{eq:features} in the chart $\Psi$}
\If{plain}
  \State $\hat g\gets\frac1{n_{\mathrm{ep}}}\sum_m\gamma^{t_m}\hat A_m\chi_m$ (blockwise);\quad $\Delta\psi\gets\mathrm{lr}\cdot\hat g$ \Comment{\eqref{eq:ghat}}
\Else
  \ForAll{visited states $s$} \Comment{\eqref{eq:local-estimates}, \eqref{eq:neff}, \eqref{eq:statewise}}
    \State compute $\hat d_s,\hat C_s,\hat b_s,n_{\mathrm{eff}}(s)$ from the visits $M_s$
    \State $\lambda_s\gets\lambda_0/\sqrt{n_{\mathrm{eff}}}$;\quad $c_s\gets n_{\mathrm{eff}}/(n_{\mathrm{eff}}+n_0)$;\quad $\hat w_s\gets c_s(\hat C_s+\lambda_sD_s^{-2})^{-1}\hat b_s$
  \EndFor
  \State $q\gets\sum_s\hat d_s\hat w_s\T\hat C_s\hat w_s$;\quad $\Delta\psi\gets\sqrt{2\varepsilon/q}\,\hat w$ \Comment{$0$ if $q\approx0$; \eqref{eq:q-alpha}}
\EndIf
\State $\Delta\psi\gets\Delta\psi\cdot\min\{1,\ c_{\max}/\|D^{-1}\Delta\psi\|_\infty\}$;\quad $\psi\gets\psi+\Delta\psi$ \Comment{\eqref{eq:cap}}
\end{algorithmic}
\end{algorithm}

The exact (noise-free) version replaces lines 2--7 of \cref{alg:update} by the closed forms \eqref{eq:grad-closed} and \eqref{eq:natural-closed}:
\begin{align*}
  \Delta\psi&=\mathrm{lr}\,D^{-1}\bigl(u_p(s)\,\pi(b\mid s)\,A(s,b)\bigr)_{s,b} &&\text{(plain)},\\
  \Delta\psi&=\sqrt{\frac{2\varepsilon}{q}}\,D\bigl(A(s,b)-A(s,\aref)\bigr)_{s,b} &&\text{(natural)},
\end{align*}
with $q=\sum_sd_p(s)\Var_\pi[A(s,\cdot)]$. It is the object that the sampled algorithm approximates.

\paragraph{Concluding remarks on Part I.}
The geometry gives three things: a direction that is a vector of $T_p\mani$ and hence covariant, a step size $\varepsilon$ with an intrinsic meaning (the Kullback--Leibler divergence, valid to second order), and, in the tabular chart,
a closed form in which neither the visitation nor the probabilities appear. The implementation needs three ingredients on top: a genuine chart (\cref{prop:chart}), a regularisation applied statewise
(\cref{sec:regularisation}), and a cap on the step. The theory does not say that the natural gradient is faster than a tuned plain gradient, nor does it address exploration; it also holds for the discrete step only to first order in the step size,
and the damping, being a numerical device, is the one place where the covariance is maintained by construction (\cref{prop:covariance-impl}) rather than by the geometry itself.

% =====================================================================
\part{Practice: the experiments}
% =====================================================================

\section{Implementation and protocol}\label{sec:protocol}

Both experiments are Jupyter notebooks in which every setting (layout, rewards, discount, batch size, grids, damping, cap, number of seeds) is a variable at the top. They apply the tabular algorithm of \cref{alg:update} to two environments.

\paragraph{Gaussian bandit.}
$K$ arms with rewards $r\sim\mathcal N(\mu_a,\sigma_a^2)$; one state, so \eqref{eq:decoupled} is a single regression. The batch is $N$ pulls, the baseline is the mean of the other rewards of the batch (independent of the pulled arm, hence unbiased), the
ridge is constant ($\lambda=\lambda_0/K$) without shrinkage ($n_0=0$), and an update is skipped for an uninformative batch.

\paragraph{Gridworld.}
An $R\times C$ grid with a start, a goal, pits and bonus cells (given or drawn at random). Bonuses are collected once per episode, so a state is $s=(c,\kappa)$ with $c$ a cell and $\kappa\in\{0,1\}^m$ the mask of the collected bonuses; $S$ is the set of non-terminal states reachable from
the start. A move into the cell $c'$ pays
\[
  r=R_{\mathrm{step}}+R_{\mathrm{goal}}\Ind[c'=\text{goal}]+R_{\mathrm{pit}}\Ind[c'\in\text{pits}]+R_{\mathrm{bonus}}\Ind[c'\text{ is an uncollected bonus}],
\]
with $R_{\mathrm{step}}$ also paid when moving against a wall; the goal and the pits end the episode. With probability $p_{\mathrm{slip}}$ the executed action is replaced by a random one (here $p_{\mathrm{slip}}=0$).

\paragraph{Exact evaluation.}
The finite MDP gives $V^\pi$, $Q^\pi$ and $u^\pi$ by linear algebra. With the transition matrix $P_\pi$ and the expected reward $r_\pi$ under $\pi$ (terminal states have value $0$),
\[
  V^\pi=(I-\gamma P_\pi)^{-1}r_\pi,\qquad
  u^\pi=(I-\gamma P_\pi\T)^{-1}e_{s_0},\qquad
  Q^\pi(s,a)=r(s,a)+\gamma\,\E\bigl[V^\pi(s')\mid s,a\bigr],
\]
and $\eta=V^\pi(s_0)$. Value iteration gives the optimum $\eta^*$, and the greedy policy of the MDP with $R_{\mathrm{bonus}}=0$, evaluated in the true MDP, gives $\eta_{\mathrm{nb}}$, the return of a policy ignoring the bonuses. These exact quantities provide the
ground truth of the checks and the noise-free runs, in which the exact gradient \eqref{eq:grad-closed} or the exact natural direction \eqref{eq:natural-closed} replaces its estimate. A run is \emph{solved} if its exact return reaches
$\eta_{\mathrm{init}}+0.95\,(\eta^*-\eta_{\mathrm{init}})$ at some iteration.

\begin{table}[!htbp]
\centering
\footnotesize
\begin{tabular}{@{}lll@{}}
\toprule
 & Bandit & Gridworld\\
\midrule
Environment & $\mu=(2,5,8)$, $\sigma^2=2$, & $5\times5$, start $(0,0)$, goal $(4,4)$,\\
 & $\pi_{\mathrm{init}}=(0.9,\,0.08,\,0.02)$ & bonuses $(0,1),(2,3),(2,4)$,\\
 & & pits $(1,1),(2,2),(3,3),(3,4),(4,1)$\\
Rewards, discount & --- & $R_{\mathrm{step}}=-1$, $R_{\mathrm{goal}}=100$, $R_{\mathrm{pit}}=-50$,\\
 & & $R_{\mathrm{bonus}}=30$, $\gamma=0.95$\\
Reference action & the last (best) arm & ``right''\\
Samples per update & $N=100$ pulls & $n_{\mathrm{ep}}=6000$ episodes, $T_{\max}=1000$\\
Regularisation & $\lambda=0.03/K=0.01$, $n_0=0$, $\tau=10^{-12}$ & $\lambda_0=0.05$, $n_0=10$\\
Cap & $c_{\max}=10$ & $c_{\max}=50$\\
Grids & $\mathrm{lr}\in\{0.1,0.5,2,8,30\}$, & $\mathrm{lr}\in\{0.003,\dots,3\}$ (7 values),\\
 & $\varepsilon\in\{0.01,0.03,0.1,0.3\}$ & $\varepsilon\in\{0.003,\dots,0.3\}$ (5 values)\\
Runs & 60 updates, 40 seeds (exact: 100 updates) & 15 updates, 10 seeds (exact: 300 updates)\\
\bottomrule
\end{tabular}
\caption{Settings used in the notebooks.}\label{tab:settings}
\end{table}

\begin{figure}[!htbp]
\centering
\figslot[0.9]{fig_bandit_rewards}
\caption{Bandit: reward distributions of the three arms and initial policy (best arm in red).}\label{fig:bandit-rewards}
\end{figure}

% ---------------------------------------------------------------------
\FloatBarrier
\section{The bandit}\label{sec:bandit}

\subsection{Setup and verification}

The best arm has mean $8$ but probability $0.02$ at the start, which gives $\eta_{\mathrm{init}}=2.360$ against $\eta^*=8$ and a success target $7.718$ (\cref{fig:bandit-rewards}). The manifold has dimension $K-1=2$.

At the initial policy and in the chart with the last arm as reference, the objects of Part~I can be written out. First,
\[
  \eta=0.9\cdot2+0.08\cdot5+0.02\cdot8=2.360 .
\]
The gradient \eqref{eq:grad-closed} and the Fisher matrix \eqref{eq:Cs} are
\[
  \nabla\tilde\eta=\bigl(\pi_b(\mu_b-\eta)\bigr)_{b=0,1}=\begin{pmatrix}-0.324\\0.2112\end{pmatrix},
  \qquad
  C=\begin{pmatrix}0.09&-0.072\\-0.072&0.0736\end{pmatrix},
\]
with $\det C=0.00144$ and eigenvalues $0.1543$ and $0.0093$, so $C$ is positive definite. The natural direction \eqref{eq:natural-closed} and its check against \eqref{eq:natural-eq} are
\[
  w=(\mu_b-\mu_{\aref})_{b=0,1}=\begin{pmatrix}-6\\-3\end{pmatrix},
  \qquad
  Cw=\begin{pmatrix}0.09\cdot(-6)-0.072\cdot(-3)\\-0.072\cdot(-6)+0.0736\cdot(-3)\end{pmatrix}=\begin{pmatrix}-0.324\\0.2112\end{pmatrix}=\nabla\tilde\eta,
\]
and the squared Fisher length is
\[
  w\T Cw=\Var_{a\sim\pi}[\mu_a]=0.9\cdot4+0.08\cdot25+0.02\cdot64-2.36^2=1.3104 .
\]
The gradient is dominated by the factors $\pi_b$ (the best arm has $\pi=0.02$ and is not even among the two coordinates, being the reference arm); the natural direction is not.

The formulas of Part~I were also checked numerically in the notebook. The exact gradient \eqref{eq:grad-closed} agrees with finite differences of $\eta$ to $2.5\cdot10^{-10}$, and the linear solution of $Fw=\bar g$ agrees with the closed form $\mu_a-\mu_{\aref}$ to $9\cdot10^{-16}$.
With a batch of $2\cdot10^5$ pulls the sampled directions have cosine $1.0000$ with the exact ones.

\paragraph{Kullback--Leibler budget.}
Taking the exact natural step of budget $\varepsilon$ and measuring the true Kullback--Leibler divergence (\cref{fig:bandit-kl}), the ratio $\mathrm{KL}/\varepsilon$ tends to $1$ as $\varepsilon\to0$ (it is $1.0$ at $10^{-4}$) and rises to about $3$ for $\varepsilon\approx1$: the budget \eqref{eq:budget} is a second-order statement, valid for small steps, and
least accurate here because the policy is close to the boundary of the simplex.

\begin{figure}[!htbp]
\centering
\figslot[0.6]{fig_bandit_kl}
\caption{Bandit: realised Kullback--Leibler divergence divided by $\varepsilon$ for the exact natural step.}\label{fig:bandit-kl}
\end{figure}

\subsection{Noise-free comparison}

With exact gradients (\cref{fig:bandit-exact}, \cref{tab:bandit-exact}) natural gradient closes the whole gap for all $\varepsilon$, and the number of iterations decreases monotonically with $\varepsilon$ (to $95\%$: $19$, $11$, $6$, $3$).
Plain gradient ascent depends strongly on the learning rate: for $\mathrm{lr}\le0.5$ it does not reach the target (the gap closed after $100$ iterations is $46\%$ for $\mathrm{lr}=0.1$ and $79\%$ for $\mathrm{lr}=0.5$), and its curves show a plateau at $\eta=5$, the value of the middle arm: the gradient
component of the best arm is proportional to $\pi=0.02$, so the policy first concentrates on the middle arm and escapes slowly. A large learning rate removes the plateau ($12$ iterations at $\mathrm{lr}=8$), but $\mathrm{lr}=30$ is slower ($30$ iterations): it overshoots onto the plateau.

\begin{figure}[!htbp]
\centering
\figslot[0.95]{fig_bandit_exact}
\caption{Bandit, exact gradients: $\eta$ against the iteration for plain gradient ascent (left) and natural gradient (right). Solid line: $\eta^*$; dashed: target.}\label{fig:bandit-exact}
\end{figure}

\begin{table}[!htbp]
\centering
\small
\begin{tabular}{@{}lrrrrr@{}}
\toprule
 & final $\eta$ & $50\%$ & $90\%$ & $95\%$ & $99\%$\\
\midrule
plain, $\mathrm{lr}=2$ & 7.986 & 23 & 29 & 31 & 45\\
plain, $\mathrm{lr}=8$ & 7.997 & 10 & 12 & 12 & 13\\
plain, $\mathrm{lr}=30$ & 7.999 & 29 & 30 & 30 & 30\\
natural, $\varepsilon=0.01$ & 8.000 & 9 & 17 & 19 & 21\\
natural, $\varepsilon=0.1$ & 8.000 & 3 & 5 & 6 & 7\\
natural, $\varepsilon=0.3$ & 8.000 & 1 & 3 & 3 & 4\\
\bottomrule
\end{tabular}
\caption{Bandit, exact gradients: iterations needed to close $50,90,95,99\%$ of the gap (excerpt).}\label{tab:bandit-exact}
\end{table}

\subsection{Sampled comparison and the role of the damping}

With samples (\cref{fig:bandit-sampled}, \cref{tab:bandit-sampled}), damped natural gradient solves all $40$ runs for every $\varepsilon$, with a worst final return of at least $7.995$, and the median number of pulls to the target decreases with $\varepsilon$ from $2150$ to $500$.
Plain gradient ascent fails for $\mathrm{lr}\le0.5$ and, for larger learning rates, ends on the plateau at $\eta\approx5$ in some runs (worst final return $5.0$): $39$, $36$ and $29$ runs out of $40$ are solved for $\mathrm{lr}=2,8,30$.

Without damping every run reaches the target but some collapse afterwards (worst final return between $3.6$ and $5.3$). The trace of the worst run ($\varepsilon=0.1$) shows the mechanism of \cref{sec:regularisation}: at iteration $7$ one arm of probability $\approx0.00$ is missing from the batch, the
fitted weights reach $\|w\|_\infty=178$, and at iteration $8$ the policy has jumped from $(0.00,\,0.018,\,0.982)$ to $(0.835,\,0.002,\,0.163)$ and $\eta$ falls from $7.945$ to $2.985$. The unsampled arm carries no information in $\hat C$, and its fitted weight is determined by round-off and by its tiny probability.

\begin{figure}[!htbp]
\centering
\figslot[0.95]{fig_bandit_sampled}
\caption{Bandit, sampled gradients ($40$ seeds): mean $\eta$ against the number of pulls for plain gradient ascent, damped natural gradient and natural gradient without damping; the last panel shows the best setting of each with its $10$--$90\%$ band.}\label{fig:bandit-sampled}
\end{figure}

\begin{table}[!htbp]
\centering
\small
\begin{tabular}{@{}lrrrr@{}}
\toprule
 & final $\eta$ & worst & solved & median pulls\\
\midrule
plain, $\mathrm{lr}=2$ & 7.892 & 5.038 & 39/40 & 3100\\
plain, $\mathrm{lr}=8$ & 7.693 & 5.001 & 36/40 & 1300\\
plain, $\mathrm{lr}=30$ & 7.175 & 5.000 & 29/40 & 1300\\
natural, $\varepsilon=0.03$ & 7.999 & 7.997 & 40/40 & 1300\\
natural, $\varepsilon=0.1$ & 8.000 & 7.997 & 40/40 & 700\\
natural, $\varepsilon=0.3$ & 8.000 & 7.998 & 40/40 & 500\\
natural, no damping, $\varepsilon=0.03$ & 7.834 & 4.386 & 40/40 & 1100\\
natural, no damping, $\varepsilon=0.1$ & 7.889 & 3.583 & 40/40 & 600\\
\bottomrule
\end{tabular}
\caption{Bandit, sampled gradients, $40$ seeds, $60$ updates of $100$ pulls (excerpt).}\label{tab:bandit-sampled}
\end{table}

\subsection{Dependence on the parametrization}

\paragraph{Rescaling.}
In a chart $\Psi=D\Theta$ with random diagonal $D$ (factors in $[1/3,3]$), the trajectories of natural gradient are unchanged up to round-off, while those of plain gradient ascent change (\cref{fig:bandit-rescale}). With the same seeds, the maximal difference in $\eta$ between the two charts is
$8.9\cdot10^{-16}$ (exact) and $2.2\cdot10^{-15}$ (sampled) for natural gradient, against $2.97$ and $2.92$ for plain gradient ascent. The sensitivity test (right panel) shows that the window of good learning rates moves by a factor of about $c^2$ when the chart is rescaled by $c=3$, as
predicted by \eqref{eq:c2}, while the curve of $\eta$ against $\varepsilon$ is the same in both charts (difference $3.6\cdot10^{-15}$).

\begin{figure}[!htbp]
\centering
\figslot[0.95]{fig_bandit_rescale}
\caption{Bandit, reparametrization. Left: exact trajectories in the original (solid) and rescaled (dashed) chart. Right: median $\eta$ after $30$ updates against the learning rate and against $\varepsilon$, for $c=1$ and $c=3$.}\label{fig:bandit-rescale}
\end{figure}

\paragraph{Reference arm.}
A change of reference arm is another chart, related to the first by the affine map \eqref{eq:change-ref}. The exact natural trajectory is the same for the three choices, while plain gradient ascent depends on it (\cref{fig:bandit-refarm}, \cref{tab:bandit-refarm}): at $\mathrm{lr}=8$ it needs $48$, $2$ or $12$ iterations
when the worst, the middle or the best arm is the reference. The reference arm changes the iteration count at a fixed learning rate by a factor of up to $24$, and changes the best learning rate ($30$ for the worst arm, $8$ for the other two). With the middle arm as reference, plain gradient ascent
reaches the target faster than natural gradient ($2$ iterations against $6$) and solves all $40$ sampled runs for $\mathrm{lr}\ge0.5$.

\begin{figure}[!htbp]
\centering
\figslot[0.95]{fig_bandit_refarm}
\caption{Bandit, exact: plain gradient ascent ($\mathrm{lr}=2$) for the three choices of reference arm (left), and natural gradient (right).}\label{fig:bandit-refarm}
\end{figure}

\begin{table}[!htbp]
\centering
\small
\begin{tabular}{@{}lrrrrrrrr@{}}
\toprule
 & \multicolumn{4}{c}{exact: iterations to target} & \multicolumn{4}{c}{sampled: solved runs of 40}\\
\cmidrule(lr){2-5}\cmidrule(lr){6-9}
reference arm & lr 2 & lr 8 & lr 30 & natural & lr 2 & lr 8 & lr 30 & natural\\
\midrule
arm 0 (mean 2) & --- & 48 & 4 & 6 & 3 & 15 & 26 & 39\\
arm 1 (mean 5) & 5 & 2 & 2 & 6 & 40 & 40 & 40 & 40\\
arm 2 (mean 8) & 31 & 12 & 30 & 6 & 39 & 36 & 29 & 40\\
\bottomrule
\end{tabular}
\caption{Bandit: dependence on the reference arm (natural gradient with $\varepsilon=0.1$; ``---'': target not reached in $100$ iterations).}\label{tab:bandit-refarm}
\end{table}

\paragraph{Important remark}
Playing with multiple configurations for the bandit setup, in particular the number of arms, the mean rewards, and the variances, one can quickly observe that in many cases the vanilla gradient actually converges faster than the natural gradient by a few iterations. Despite this fact, this experiment demonstates multiple key advantages of the natural gradient over the vanilla gradient. First, the choice of initial policy which places very little weight on the best arm is an example of the sort of challenging scenario an agent might face in a more general MDP: the best action may not have the highest probability initially, and even with uniform initial distribution visitation biases can develop after a few rounds.
Since in the vanilla gradient the update is proportional to the visitation, it can be very slow to escape from a local optimum. The natural gradient, on the other hand, cancels dependence on the visitation through the inverse of the Fisher matrix, and can thus escape from such plateauing behaviour -- this phenomenon is clearly reflected in the experiment above. 
Second, the natural gradient was shown to be invariant under reparametrization of the policy, which is not true for the vanilla gradient. The experiment shows that arbitrary parametrization decisions, which are purely mathematical artefacts and should not impact the convergence of the algorithm, can have a significant effect on the convergence of the vanilla gradient and the expected return at its converged value. 
Third, one could make the argument that the natural gradient is more rubust to hyperparameter tuning, as it converges at acceptable rates for a wide range of $\varepsilon$ values (from 0.01 to 0.3), while the vanilla gradient is very sensitive to the choice of learning rate, and values of $\alpha$ must be tested over multiple orders of magnitude (0.1 -> 1 -> 10) in order to find a good value. 
Finally, the natural gradient is more robust to noise in the gradient estimate, as it was shown that the natural gradient converged to the target in all 40 runs for all $\varepsilon$ values, while the vanilla gradient failed to converge in many runs for all learning rates except for lr=2 with the middle arm as reference.

% ---------------------------------------------------------------------
\FloatBarrier
\section{The gridworld}\label{sec:grid}

\subsection{Setup and verification}

The layout of \cref{fig:grid-layout} has $|S|=82$ reachable states, so $\dim\mani=82\cdot3=246$. The uniform policy returns $\eta_{\mathrm{init}}=-22.9$, the optimum is $\eta^*=102.55$, and the best policy ignoring the bonuses returns $\eta_{\mathrm{nb}}=63.10$; the optimal route collects the three bonuses before
reaching the goal. The target is $-22.9+0.95\,(102.55+22.9)=96.28$.

The checks of Part~I hold: the exact gradient agrees with finite differences to $1.0\cdot10^{-8}$, and the solution of $Fw=\bar g$ agrees with the closed form $A(s,a)-A(s,\aref)$ to $1.1\cdot10^{-13}$. The Fisher matrix is positive definite, but its eigenvalues range from $2.0\cdot10^{-12}$
to $8.1\cdot10^{-2}$ at a random policy: it is numerically near-singular, because the weight $d_p(s)$ of rarely visited states is tiny. This is the situation described in \cref{rem:damping}. The maps of \cref{fig:grid-layout} illustrate \cref{prop:natural-closed}: the plain-gradient update is concentrated
near the start, where $d_p$ is large (largest component about $7$), whereas the natural update is not tied to the visitation: its largest component, above $200$, is in the cell next to the goal.

\begin{figure}[!htbp]
\centering
\figslot[0.95]{fig_grid_layout}
\caption{Gridworld. Left: layout (bonuses B0--B2, pits, goal). Right: at the initial policy and for the states without collected bonus, visitation $d_p(s)$, size of the plain-gradient update and size of the natural update in each cell.}\label{fig:grid-layout}
\end{figure}

\subsection{Noise-free comparison}

With exact gradients (\cref{fig:grid-exact}, \cref{tab:grid-exact}), natural gradient reaches $\eta^*$ ($100\%$ of the gap) for every $\varepsilon$, monotonically, and reaches the target in $25$ to $4$ iterations as $\varepsilon$ grows from $0.003$ to $0.3$. Plain gradient ascent never exceeds $64.2\%$ of the gap: for every learning rate
it ends below $58$, a policy that does not collect the bonuses ($57.6<\eta_{\mathrm{nb}}$), and the gradient there is small because it carries the factors $u_p(s)\pi(a\mid s)$. Larger learning rates only reach this level faster.

\begin{figure}[!htbp]
\centering
\figslot[0.95]{fig_grid_exact}
\caption{Gridworld, exact gradients: plain gradient ascent (left) and natural gradient (right). Solid: $\eta^*$; dotted: $\eta_{\mathrm{nb}}$; dashed: target.}\label{fig:grid-exact}
\end{figure}

\begin{table}[!htbp]
\centering
\small
\begin{tabular}{@{}lrrrrrr@{}}
\toprule
 & final $\eta$ & gap closed & $50\%$ & $90\%$ & $95\%$ & $99\%$\\
\midrule
plain, $\mathrm{lr}=0.03$ & 52.86 & 60.4\% & 109 & --- & --- & ---\\
plain, $\mathrm{lr}=0.3$ & 57.48 & 64.1\% & 12 & --- & --- & ---\\
plain, $\mathrm{lr}=3$ & 57.63 & 64.2\% & 3 & --- & --- & ---\\
natural, $\varepsilon=0.003$ & 102.55 & 100\% & 10 & 22 & 25 & 29\\
natural, $\varepsilon=0.03$ & 102.55 & 100\% & 4 & 8 & 9 & 11\\
natural, $\varepsilon=0.3$ & 102.55 & 100\% & 2 & 4 & 4 & 5\\
\bottomrule
\end{tabular}
\caption{Gridworld, exact gradients: iterations (out of $300$) needed to close a fraction of the gap (excerpt).}\label{tab:grid-exact}
\end{table}

\subsection{The sampled estimator and covariance}

With a batch of $4000$ episodes ($37512$ steps) at the initial policy, $62$ of the $82$ states are visited. The estimate of the plain gradient is accurate (cosine $0.999$ with the exact gradient, relative error $0.037$), because it is dominated by the frequently visited states. The estimate of the natural direction is much noisier
(cosine $0.45$, relative error $0.90$): it gives all states the same weight, and the effective number of visits per visited state ranges from $1.0$ (minimum) through $54.1$ (median) to $9884$ (maximum). The noise sits in the rarely visited states, where the shrinkage $c_s$ and the ridge $\lambda_s$ of \eqref{eq:statewise} act;
its effect on the trajectories is measured below.

Under a rescaling of the chart with random $D$, the natural-gradient trajectories change by less than $10^{-12}$ in $\eta$ (exact and sampled; round-off level), whereas plain gradient ascent changes by $11.1$ in both cases (\cref{fig:grid-rescale}); this confirms \cref{prop:covariance-impl}, including the statewise ridge $\lambda_sD_s^{-2}$.

\begin{figure}[!htbp]
\centering
\figslot[0.7]{fig_grid_rescale}
\caption{Gridworld: trajectories in the original chart (solid) and in the rescaled chart (dashed), exact gradients.}\label{fig:grid-rescale}
\end{figure}

\subsection{Sampled comparison}

\Cref{fig:grid-sampled} and \cref{tab:grid-sampled} give the runs with samples ($10$ seeds, $15$ updates of $6000$ episodes). Plain gradient ascent solves none of the runs: its best setting reaches $57.4$ after $15$ updates, which is the exact plateau ($57.6$). Natural gradient with $\varepsilon\ge0.03$ solves $8$, $8$ and $7$ of the $10$ runs,
after a median of $69000$, $45000$ and $36000$ episodes; with $\varepsilon=0.003$ and $0.01$ the budget is too small for $15$ updates (final returns $63.8$ and $86.5$). The sampled curves follow the exact ones with a delay: the exact curves reach the target after $9$, $6$ and $4$ iterations, i.e.\ $54000$, $36000$ and $24000$ episodes,
against the sampled medians above, an additional $1.5$ to $2$ iterations. This shows that the statewise regularisation recovers the behaviour of the exact natural gradient despite the noisy direction estimate.

The runs that are not solved end on a route with a single bonus and a return of exactly $85.00$: the policy has saturated on a one-bonus route before the other bonuses have been explored, so that exploration, not the geometry of the update, is the limiting factor. This accounts for the gap between the mean final return ($99.0$) and $\eta^*$:
for instance $(8\cdot102.55+2\cdot85)/10=99.04$.

\begin{figure}[!htbp]
\centering
\figslot[0.95]{fig_grid_sampled}
\caption{Gridworld, sampled gradients ($10$ seeds): mean $\eta$ against the number of episodes for plain gradient ascent (left) and natural gradient (centre); right: best setting of each with its $10$--$90\%$ band.}\label{fig:grid-sampled}
\end{figure}

\begin{table}[!htbp]
\centering
\small
\begin{tabular}{@{}lrrrrl@{}}
\toprule
 & final $\eta$ & worst & solved & median episodes & final routes (bonuses)\\
\midrule
plain, $\mathrm{lr}=0.3$ & 46.53 & 46.38 & 0/10 & --- & goal not reached (10)\\
plain, $\mathrm{lr}=3$ & 57.38 & 57.36 & 0/10 & --- & goal not reached (10)\\
natural, $\varepsilon=0.01$ & 86.52 & 84.21 & 0/10 & --- & 1 bonus (8), 3 bonuses (2)\\
natural, $\varepsilon=0.03$ & 99.02 & 85.00 & 8/10 & 69000 & 1 bonus (2), 3 bonuses (8)\\
natural, $\varepsilon=0.1$ & 99.04 & 85.00 & 8/10 & 45000 & 1 bonus (2), 3 bonuses (8)\\
natural, $\varepsilon=0.3$ & 97.29 & 85.00 & 7/10 & 36000 & 1 bonus (3), 3 bonuses (7)\\
\bottomrule
\end{tabular}
\caption{Gridworld, sampled gradients: final return, worst seed, number of solved runs, median number of episodes to the target, and bonuses collected by the final greedy route (excerpt; the target is $96.28$).}\label{tab:grid-sampled}
\end{table}

% =====================================================================
\FloatBarrier
\section{Conclusion}\label{sec:conclusion}

Policies form a manifold, and the Fisher metric, the Hessian of the Kullback--Leibler divergence at its minimum, is an intrinsic metric on it: with it the covariant gradient becomes a contravariant direction, and the step size $\varepsilon$ is a Kullback--Leibler budget. In the reference-logit chart the metric is block diagonal and positive definite,
and the natural gradient is the advantage relative to the reference action, in which neither the visitation nor the policy probabilities appear.

The experiments confirm this on two environments. The natural-gradient trajectories are independent of the chart to round-off, whereas those of plain gradient ascent change with the scale of the coordinates and with the reference action. In the gridworld, exact natural gradient reaches the optimum for every $\varepsilon$, while plain gradient ascent
stalls on a policy that ignores the bonuses, whose gradient carries the small factors $u_p(s)\pi(a\mid s)$. The sampled version needs a regularisation that respects the structure: applied to the global system the ridge reintroduces the visitation (\cref{rem:global-ridge}); applied statewise it lets the sampled trajectories follow the exact ones within about two iterations.

The limitations are the exact baseline $V^\pi$, tabular policies, short runs ($10$ seeds and $15$ updates in the gridworld), a noisy natural direction in rarely visited states (cosine $0.45$ at $4000$ episodes), and convergence of some runs to a one-bonus route, a matter of exploration that the natural gradient does not address. In the bandit, a tuned plain gradient in a favourable chart is as fast as
natural gradient: the advantage of the latter lies in invariance, in a step size with an intrinsic meaning and in robustness to the learning rate, not in speed in every chart. Natural extensions are a learned critic, function approximation (where the regression of \cref{thm:kakade} is no longer exact), and geodesic steps for the Fisher metric.

\clearpage

\begin{thebibliography}{9}
\bibitem{kakade2001} S.~Kakade. A natural policy gradient. \emph{Advances in Neural Information Processing Systems} 14, 2001.
\bibitem{sutton2000} R.~S.~Sutton, D.~McAllester, S.~Singh, Y.~Mansour. Policy gradient methods for reinforcement learning with function approximation. \emph{Advances in Neural Information Processing Systems} 12, 2000.
\bibitem{sutton2018} R.~S.~Sutton, A.~G.~Barto. \emph{Reinforcement Learning: An Introduction}, 2nd ed., MIT Press, 2018.
\bibitem{williams1992} R.~J.~Williams. Simple statistical gradient-following algorithms for connectionist reinforcement learning. \emph{Machine Learning} 8:229--256, 1992.
\bibitem{gibbs} Gibbs' inequality. Wikipedia, \url{https://en.wikipedia.org/wiki/Gibbs'_inequality}.
\end{thebibliography}

\end{document}
